\section{Structure of the proof}
\label{ch:proof-structure}

This chapter records the main transformations of~\cite{JNVWY20} at headline level: for
each, what it does, what it relies on, and one theorem statement so that the dependency
graph reflects the architecture. Working out each section --- restating it in the
synchronous framework, decomposing it into lemmas, and setting the constants --- is a
main task of the project. The paper factors several of these transformations through
\emph{typed} verifiers and a detyping step; that machinery is not
reproduced here and should be introduced locally where needed.

Following Lin~\cite{Lin25}, every soundness statement below is in \emph{value form}: a
transformation preserves the property ``value at most $1 - \eps$'' up to a quantitative
loss, and compression preserves ``value at most $\frac12$''. The three inner
transformations, repetition and compression all carry that property in $\valstar$, the
bipartite tensor-product value: it is the class the compressibility criterion is
instantiated with and the form Definition~\ref{def:mipstar} requires
(Remark~\ref{rem:compression-abstract}), and no passage between the two values occurs
anywhere on the path. The synchronous value is not used, and neither is the
almost-synchronous transport of Theorem~\ref{thm:almost-sync}; see
Remark~\ref{rem:sync-invariant} for why that is the better arrangement and
Remark~\ref{rem:bipartite-route} for the evidence. No entanglement lower bound is tracked. Remark~\ref{rem:direct-vs-anchored} explains why none is needed,
Section~\ref{sec:rr-compression} states the abstract lemma (Lin's compressibility
criterion) that turns value-form compression into the halting reduction, and the
entanglement clauses of~\cite{JNVWY20} are recorded for reference in
Remark~\ref{rem:entanglement-form}.

\subsection{Conditionally linear functions and samplers}
\label{sec:ps-cl}

Question distributions must be structured enough that players can sample them
\emph{on their own} in a certifiable way (introspection), yet expressive enough to
implement all the verifiers in the pipeline. The right class is that of conditionally
linear distributions: linear maps applied in stages, where each stage may depend on
the output of the previous ones.

\begin{definition}[Conditionally linear function]
\label{def:cl-function}
Let $V = \F_q^s$ and $\ell \geq 1$. An \emph{$\ell$-level conditionally linear (CL)
function} on $V$ is defined recursively: a $1$-level CL function is a linear map $L :
V \to V$; an $\ell$-level CL function is given by a direct sum decomposition $V = V_1
\oplus V_{>1}$, a linear map $L_1$ on $V_1$, and, for each value $u$ of $L_1$, an
$(\ell-1)$-level CL function $L^{u}_{>1}$ on $V_{>1}$; it maps $z = z_1 + z_{>1}$ to
$L_1(z_1) + L^{L_1(z_1)}_{>1}(z_{>1})$. A \emph{CL distribution} is the distribution of
$(L^\alice(z), L^\bob(z))$ for uniformly random $z \in V$ and CL functions $L^\alice,
L^\bob$.
\end{definition}

An $\ell$-level sampler (Definition~\ref{def:sampler}) presents a CL distribution via
Turing machines answering structured queries (dimensions, marginals, and evaluations of
the linear maps and their canonical complements). This section of the project should
fix the query interface once and for all, and prove the closure properties used later:
products, concatenation, and reduction of the field size to $q = 2$
(``downsizing'', using Lemma~\ref{lem:self-dual-basis}).

\begin{remark}[Why the pipeline is bipartite throughout]
\label{rem:bipartite-route}
\uses{def:q-value,def:sync-value,thm:almost-sync,rem:almost-sync-hypothesis,rem:sync-invariant,lem:sync-le-valstar}
Every soundness clause of this chapter is stated in $\valstar$ and none in $\synval$. The
choice is deliberate and rests on two things.

First, it is forced at the outer end. Applied to a game supported on the diagonal the
corrected transport (Theorem~\ref{thm:almost-sync}) forces its constant to satisfy $C \geq
2^{c-1}$, so the best it can conclude from $\synval \leq \frac12$ is $\valstar \leq 1 -
\kappa (2C)^{-1/c}$, which is never better than $\frac12$; and its diagonal-weight
hypothesis is in any case unavailable on a $k$-fold product, whose diagonal weight decays
as $\kappa^{k}$ (Remark~\ref{rem:almost-sync-hypothesis}). So the threshold has to come
from a $\valstar$ theorem, which Theorem~\ref{thm:direct-repetition-q} is.

Second, the inner end does not need $\synval$ either. The adversarially verified
decomposition of the proof recorded in the companion repository
\texttt{vidick/mipre-proof} --- 123 nodes, 291 challenges all resolved --- states every
node in the bipartite tensor-product value, from the top claim to the conclusion; not one
of its nodes mentions the synchronous value. Stating this chapter's clauses in $\synval$
therefore buys nothing the sources provide and costs three things: the transport itself
(Theorem~\ref{thm:almost-sync}, whose printed form was false), its polynomial loss and the
factor $\kappa^{13}$ it contributed to
Theorem~\ref{thm:parallel-repetition}, and the invariant of
Remark~\ref{rem:sync-invariant}, which had to be established for three transformations and
never was. All three are now off the path. Lemma~\ref{lem:sync-le-valstar} survives as the
free comparison it always was, used only to read a $\valstar$ conclusion in $\synval$ when
a synchronous statement is wanted (Theorem~\ref{thm:halting}, Theorem~\ref{thm:main}),
never in the other direction.

What this costs in turn is that the rigidity and approximate-measurement lemmas the inner
analyses consume must be stated bipartitely, as their sources state them, rather than
tracially. That is the arrangement of~\cite{JNVWY20} and of the companion repository's
paper, so it is the arrangement with the least to prove.
\end{remark}

\subsection{Question reduction: introspection}
\label{sec:ps-introspection}

Introspection removes the verifier's cost of sampling questions: instead of receiving
questions for $\verifier_{2^n}$, the players \emph{sample the questions themselves} by
measuring shared EPR pairs, and the new verifier uses the Pauli basis test
(Theorem~\ref{thm:qld}) to force honest sampling. The CL structure of the sampler is
what makes commitment to each stage of the sampled question enforceable (the ``hiding''
directions of the CL decomposition). The result is a verifier whose questions are
exponentially smaller, at the price of a fixed number of levels.

\begin{theorem}[Introspection]
\label{thm:introspection}
\uses{def:normal-verifier,def:lambda-bounded,def:pcc,thm:qld,def:cl-function}
There is a polynomial-time Turing machine $\ComputeIntroVerifier$ that on input
$(\verifier, \lambda, \ell)$ returns a $5$-level normal form verifier $\vint =
(\sint, \dint)$ with
$\TIME_{\sint}(n) = \poly(n, \lambda, \ell)$,
$\TIME_{\dint}(n) = \poly(2^{\lambda n}, \ell)$ and $\abs{\dint} = \poly(\lambda,
\ell)$, such that for every $\lambda$-bounded $\ell$-level verifier $\verifier$ there
is a function $\delta(\eps, n) = a\bigl( (\lambda n)^a \eps^b + (\lambda n)^{-b}
\bigr)$ (constants $a, b$ depending only on $\ell$) with, for all $n$:
\begin{enumerate}
\item (\textbf{Completeness}) If $\verifier_{2^n}$ has a value-$1$ PCC strategy, so
  does $\vint_n$.
\item (\textbf{Soundness}) If $\valstar(\vint_n) > 1 - \eps$ then
  $\valstar(\verifier_{2^n}) \geq 1 - \delta(\eps, n)$.
\end{enumerate}
\end{theorem}

The soundness analysis is a long rigidity argument~\cite[Section 8]{JNVWY20}. It builds on techniques introduced in~\cite{NW19}, and it may be advisable to formalize some of the key lemmas there first. These include the Pauli basis test together with a toolkit of approximate-measurement lemmas
(twirling, commutation, sandwiching). For this project, those lemmas should be stated and proved in the bipartite setting, as in the sources: the pipeline carries $\valstar$ throughout, so no tracial restatement is needed and none should be attempted (Remark~\ref{rem:bipartite-route}).

\subsection{Oracularization}
\label{sec:ps-oracularization}

Oracularization converts a normal form verifier into one where a single player (the
``oracle'') receives both original questions and answers both, while the other player
receives one of the two questions and must answer consistently. This makes the game's
predicate a pointwise function of one player's answer --- the format required by answer
reduction. This is probably the easiest transformation in the paper and could be attempted first to fix the semantics. 

\begin{theorem}[Oracularization]
\label{thm:oracularization}
\uses{def:normal-verifier,def:pcc}
There is a polynomial-time Turing machine $\ComputeOracleVerifier$ mapping a normal
form verifier $\verifier = (\sampler, \decider)$ to a normal form verifier $\vora$
such that, for some $\delta(\eps) = \poly(\eps)$ and all $n$:
\begin{enumerate}
\item (\textbf{Completeness}) If $\verifier_n$ has a value-$1$ PCC strategy then
  $\vora_n$ has a value-$1$ PCC strategy.
\item (\textbf{Soundness}) If $\valstar(\vora_n) > 1 - \eps$ then
  $\valstar(\verifier_n) \geq 1 - \delta(\eps)$.
\item (\textbf{Complexity}) The sampler of $\vora$ depends only on $\sampler$, with
  $\TIME$ within a constant factor and the same number of levels.
\end{enumerate}
\end{theorem}

\subsection{Answer reduction}
\label{sec:ps-answer-reduction}

After introspection the decider still runs in time exponential in $n$ (it simulates
$\decider$ at index $2^n$), and the players' answers are correspondingly long. Answer
reduction composes the game with a PCP: by Theorem~\ref{thm:succinct-sat} the
statement ``the original decider would have accepted our answers'' is a succinct 3SAT
instance; the players commit to a low-degree encoding of a satisfying assignment
(their original answers plus the computation tableau) and the new decider spot-checks
it, using the low individual degree test to certify the encoding and
Lemma~\ref{lem:schwartz-zippel} for the arithmetization. Answers and decider time
become polylogarithmic in the original.

\begin{theorem}[Answer reduction]
\label{thm:answer-reduction}
\uses{def:normal-verifier,def:pcc,thm:succinct-sat,lem:schwartz-zippel,thm:qld,lem:self-dual-basis}
There is a polynomial-time Turing machine $\ComputeAnsVerifier$ that on input
$(\verifier, \lambda, \ell, \mu, \sigma)$, for $\verifier$ an $\ell$-level normal form
verifier with $\abs{\decider} \leq \sigma$, $\TIME_\sampler(n) \leq (\lambda n)^\mu$
and $\TIME_\decider(n) \leq (2^{\lambda n})^\mu$ for $n \geq 2$, outputs a
$\max\{\ell + 2, 5\}$-level normal form verifier $\vans$ with
$\TIME_{\sans}(n), \TIME_{\dans}(n) = \poly\bigl( (\lambda n)^\mu, \sigma \bigr)$,
whose sampler depends only on $(\sampler, \lambda, \ell, \mu, \sigma)$, and such that for
some $\delta(\eps, n) = \sigma^a \bigl( (\lambda n)^{\mu a} \eps^b + (\lambda
n)^{-\mu b} \bigr)$ (universal constants $a, b$) and all $n$:
\begin{enumerate}
\item (\textbf{Completeness}) If $\verifier_n$ has a value-$1$ PCC strategy
  then $\vans_n$ has a value-$1$ PCC strategy.
\item (\textbf{Soundness}) If $\valstar(\vans_n) > 1 - \eps$ then
  $\valstar(\verifier_n) \geq 1 - \delta(\eps, n)$.
\end{enumerate}
\end{theorem}

This is the most technically heterogeneous section of the paper
(\cite[Section 10]{JNVWY20}): half classical (tableaux, Tseitin encoding,
arithmetization --- see Section~\ref{sec:rr-cook-levin}) and half quantum (a
low-degree ``sandwiching'' argument recombining the players' committed polynomials). It may also be the hardest. 

\subsection{Gap amplification: parallel repetition}
\label{sec:ps-parallel-repetition}

Answer reduction leaves a shrunken gap ($1$ vs.\ $1 - 1/\poly(n)$). Direct parallel
repetition (Theorem~\ref{thm:direct-repetition-q}, formalized end to end) restores
soundness
$\frac12$ with polynomially many repetitions, at polynomial cost in complexity and without
anchoring; the anchored, entanglement-form theorem used in~\cite{JNVWY20}
(Theorem~\ref{thm:bvy}) is not needed (Remark~\ref{rem:direct-vs-anchored}).

\begin{theorem}[Parallel repetition of normal form verifiers]
\label{thm:parallel-repetition}
\uses{def:normal-verifier,def:pcc,def:direct-repetition,thm:direct-repetition-q,def:mipstar}
There are universal constants $b, c > 0$ such that the following holds for every
$\kappa \in (0,1]$. There are a constant $c' > 0$, depending only on $\kappa$, and a
polynomial-time Turing machine $\ComputeParrepVerifier$ that on input
$(\verifier, \lambda, \ell, \tau)$, for $\verifier$ an $\ell$-level normal form verifier,
outputs an $\ell$-level normal form verifier $\vrep$ whose game $\vrep_n$ is the
$k(n)$-fold direct repetition (Definition~\ref{def:direct-repetition}) of $\verifier_n$
--- up to the encoding of answer tuples as strings, the decider parsing a self-delimiting
$k(n)$-tuple and rejecting malformed answers --- with $k(n) = (\lambda n)^{(1 + c')\tau}$,
$\TIME_{\srep}(n) = O(k(n) \TIME_\sampler(n))$ and $\TIME_{\drep}(n) = O\bigl(k(n)
\max(\TIME_\decider(n), (\lambda n)^\tau)\bigr)$, the sampler of $\vrep$ depending only on
$(\sampler, \lambda, \ell, \tau)$, and for all $n$:
\begin{enumerate}
\item (\textbf{Completeness}) If $\verifier_n$ has a value-$1$ PCC strategy and
  $\TIME_\decider(n) \leq (\lambda n)^\tau$, then $\vrep_n$ has a value-$1$ PCC
  strategy.
\item (\textbf{Soundness}) For all $\eps > 0$, if $\valstar(\verifier_n) \leq 1 - \eps$
  and $\TIME_\decider(n) \leq (\lambda n)^\tau$, then
  $\valstar(\vrep_n) \leq \exp\bigl( -c\, \eps^{13}\, k(n) / (\lambda n)^{\tau} \bigr)$;
  in particular $\valstar(\vrep_n) \leq \frac12$ whenever $\eps \geq (\lambda
  n)^{-\tau}$, for $c'$ large enough and $n \geq 2$.
\end{enumerate}
\end{theorem}

The game $\vrep_n$ is $\verifier_n^{k(n)}$: its sampler runs $k(n)$ independent copies of
$\sampler$ (a product of conditionally linear functions of the same level) and its
decider runs $\decider$ on every coordinate. Completeness: the $k(n)$-fold tensor power of
a value-$1$ PCC strategy is a value-$1$ PCC strategy for the repeated game. Soundness is
now a direct application of Theorem~\ref{thm:direct-repetition-q}: the hypothesis is
already a $\valstar$ bound, so nothing has to be transported first. That theorem bounds
$\valstar(\verifier_n^{k(n)})$ by $\exp\bigl( -c\, k(n)\, \eps^{13} / (\eps +
\log(\abs{\mA}\abs{\mB})) \bigr)$, and the answer alphabets of $\verifier_n$ have size at
most $2^{\TIME_\decider(n)}$, so $\log(\abs{\mA}\abs{\mB}) \leq 2 (\lambda n)^\tau$ and
the denominator is $O((\lambda n)^\tau)$; with $\eps \geq (\lambda n)^{-\tau}$ and $k(n) =
(\lambda n)^{(1 + c')\tau}$ the exponent is at least $(\lambda n)^{(c' - 13)\tau}$, so
$c' > 13$ suffices. The exponent $13$ is the one of
Theorem~\ref{thm:direct-repetition-q} itself, with no further loss: carrying the
hypothesis in $\valstar$ rather than $\synval$ removes both the transport's polynomial
loss and the factor $\kappa^{13}$ that its diagonal-weight hypothesis used to contribute.
The conclusion is in $\valstar$, as is the hypothesis: that is the form
Definition~\ref{def:mipstar} asks of the final game, and it is available at every step
because nothing in the chain ever leaves $\valstar$ (Remark~\ref{rem:bipartite-route}).
The exponent $b$ and the constants are provisional (Section~\ref{sec:conventions}).

\subsection{The compression theorem}
\label{sec:ps-compression}

Composing the four transformations --- introspection, oracularization, answer
reduction, repetition --- yields gap-preserving compression: the paper-level form of
the compression hypothesis of Lemma~\ref{lem:compressible-criterion}. The soundness
clause is the value-form clause that~\cite{JNVWY20} state first and then strengthen to
an entanglement bound (Remark~\ref{rem:entanglement-form}).

\begin{theorem}[Gap-preserving compression]
\label{thm:compression}
\ledgernode{1.5, 1.5.7}
\uses{thm:introspection,thm:oracularization,thm:answer-reduction,thm:parallel-repetition,def:lambda-bounded,def:pcc}
There is a universal constant $C_0 > 0$ and a polynomial-time Turing machine
$\Compress$ that on input $(\verifier, \lambda)$ outputs a $7$-level normal form
verifier $\vcompr = (\scompr, \dcompr)$ with $\TIME_{\scompr}(n),
\TIME_{\dcompr}(n) = \poly(n, \lambda)$, where $\scompr$ is independent of
$\verifier$ and computable from $\lambda$ in time $\polylog(\lambda)$. If moreover
$\verifier$ is a $\lambda$-bounded $7$-level normal form verifier, then for all $n
\geq C_0$ and $N = 2^n$:
\begin{enumerate}
\item (\textbf{Completeness}) If $\verifier_N$ has a value-$1$ PCC strategy, then so
  does $\vcompr_n$.
\item (\textbf{Soundness}) If $\valstar(\verifier_N) \leq \tfrac12$ then
  $\valstar(\vcompr_n) \leq \tfrac12$: the class ``$\valstar \leq \tfrac12$'' is
  preserved, in the same value at both ends.
\end{enumerate}
\end{theorem}

The level count is forced by the chain rather than chosen: introspection outputs a
$5$-level verifier whatever the level of its input (Theorem~\ref{thm:introspection}),
oracularization preserves the level (Theorem~\ref{thm:oracularization}), answer reduction
takes $\ell$ to $\max\{\ell + 2, 5\}$, here $7$ (Theorem~\ref{thm:answer-reduction}), and
direct repetition preserves it (Theorem~\ref{thm:parallel-repetition}); so $\Compress$
outputs a $7$-level verifier on every input, and $7$ is consequently the level at which
the recursion closes. A verifier of fewer levels is presented at that level by padding its
decomposition with trivial factors, which an $\ell$-level conditionally linear function
(Definition~\ref{def:cl-function}) admits. In~\cite{JNVWY20} the count is $9$, anchored
repetition adding two levels (Remark~\ref{rem:direct-vs-anchored}).

The four stages are chained through a parameter budget rather than composed blindly:
each stage's soundness clause is applied contrapositively, and each application needs a
strict margin between the loss it suffers and the gap it is handed. Those margins are
arithmetic facts about the imported constants, isolated here as
Lemma~\ref{lem:compress-margin} and Lemma~\ref{lem:compress-tau}. The provenance notes
name the corresponding nodes of the companion repository's ledger
(\texttt{vidick/mipre-proof}); see \texttt{planning/ledger-informed-plan.md}.

\begin{lemma}[The compressed sampler does not depend on the verifier]
\label{lem:compress-sampler-indep}
\ledgernode{1.5.1}
\uses{thm:compression,def:normal-verifier,def:cl-function}
For all inputs $(\verifier, \lambda)$ --- $\lambda$-bounded or not, well-formed or not
--- the sampler $\scompr$ of $\Compress(\verifier, \lambda)$ depends only on $\lambda$,
and there is a Turing machine $\mathtt{ComputeSampler}$ that outputs its description from
$\lambda$ in time $\polylog(\lambda)$.
\end{lemma}
\begin{proof}
Chain the verifier-independence of the sampler through the three stages: introspection's
sampler depends only on $(\lambda, \ell)$, answer reduction's only on its input sampler
and the parameters, and repetition's only on its input sampler and $(\lambda, \ell,
\tau)$. $\mathtt{ComputeSampler}$ is then realized by running $\Compress$ on the trivial
verifier. It is this lemma that makes the recursion of Section~\ref{sec:ps-halting}
possible at all: a fixed point has to quote its own sampler, so the sampler must not
depend on what is being compressed.
\end{proof}

\effortMedium

\begin{lemma}[The margin claim]
\label{lem:compress-margin}
\ledgernode{1.5.3}
\uses{thm:introspection,thm:answer-reduction}
Write $(a_1, b_1)$ for the constants of Theorem~\ref{thm:introspection} and $(a_2, b_2)$
for those of Theorem~\ref{thm:answer-reduction}, normalized so that $a_1, a_2 \geq 1$.
Put $\sigma = \lceil C_{\mathrm{intro}} \lambda^{C_{\mathrm{intro}}} \rceil$ and
\begin{equation*}
  \mu = \Bigl\lceil \max\Bigl\{ C_{\mathrm{intro}},\ \frac{9 a_1 + 2 a_2
  C_{\mathrm{intro}}}{b_1 b_2} \Bigr\} \Bigr\rceil .
\end{equation*}
Then $\mu$ is a universal constant, $\mu \geq C_{\mathrm{intro}}$, and
\begin{equation*}
  (\lambda n)^{-b_2 \mu} \;<\; \frac{\eps_1}{2 \sigma^{a_2}}
  \qquad \text{for all } n \geq 2 \text{ and all } \lambda \geq 1,
\end{equation*}
where $\eps_1 = \bigl( 8 a_1 (\lambda n)^{a_1} \bigr)^{-1/b_1}$.
\end{lemma}

\effortEasy

\begin{lemma}[A universal repetition exponent]
\label{lem:compress-tau}
\ledgernode{1.5.5}
\uses{thm:parallel-repetition,thm:answer-reduction}
There is a universal integer $\tau$ such that, with $k(n) = (\lambda n)^{\tau}$ and
$\eps_2$ the function of $(n, \lambda)$ fixed by the answer-reduction stage, the
repetition bound of Theorem~\ref{thm:parallel-repetition} is below $\frac12$ for all $n
\geq \tau$ and all integers $\lambda \geq 1$.
\end{lemma}
\begin{proof}
$1/\eps_2$ is bounded by a fixed polynomial in $\lambda n$, because $\eps_1$, $\sigma$
and $\mu$ are, so $\eps_2^{13} \geq (\lambda n)^{-P}$ for a universal $P$; choosing $\tau$
larger than $P$ plus the answer-reduction exponent, with slack for the constants, makes
the exponent of the bound grow as a positive fixed power of $\lambda n$. Purely
arithmetic, and computationally testable --- the companion repository tests it.
\end{proof}

\effortEasy

\begin{remark}[How the value chain composes, and what it avoids]
\label{rem:compression-chain}
\ledgernode{1.5.8, 1.5.2, 1.5.4, 1.5.6}
\uses{thm:compression,thm:introspection,thm:oracularization,thm:answer-reduction,thm:parallel-repetition,def:ent,lem:compress-margin,lem:compress-tau}
The soundness clause of Theorem~\ref{thm:compression} is obtained by applying each
transformation's soundness clause contrapositively, with a margin at every step. From
$\valstar(\verifier_N) \leq \frac12$ and $\delta_1(\eps_1, n) < \frac12$: were
$\valstar(\verifier^{(1)}_n) > 1 - \eps_1$, Theorem~\ref{thm:introspection} would give
$\valstar(\verifier_N) \geq 1 - \delta_1 > \frac12$, so
$\valstar(\verifier^{(1)}_n) \leq 1 - \eps_1$. The same step through
Theorem~\ref{thm:oracularization} and then Theorem~\ref{thm:answer-reduction}, using
$\delta_2(\eps_2, n) < \eps_1$ (Lemma~\ref{lem:compress-margin}), gives
$\valstar(\verifier^{(2)}_n) \leq 1 - \eps_2$; and
Theorem~\ref{thm:parallel-repetition} with Lemma~\ref{lem:compress-tau} brings that below
$\frac12$.

Every step of that chain is an implication between $\valstar$ bounds, which is worth
recording because the route of~\cite{JNVWY20} cannot take it. There the repetition step's
soundness hypothesis is an \emph{entanglement} bound, and $\Ent(\game, 1 - \eps) = \infty$
needs $\valstar(\game) < 1 - \eps$ \emph{strictly} --- a non-strict $\leq$ leaves the value
possibly attained and the entanglement requirement finite. So the paper runs the whole
chain in $\Ent$, starting it at the threshold $1 - \delta_1$ rather than at $\frac12$; the
value-form route would force it to apply repetition at $\eps_2/2$ instead of $\eps_2$, at a
cost of $2^{17}$ in the exponent of its repetition parameter. Direct repetition
(Theorem~\ref{thm:direct-repetition-q}) takes a value hypothesis and so needs no
strictness and pays no such factor: this is the second place, after
Remark~\ref{rem:direct-vs-anchored}, where the value-form pipeline is cheaper than the
entanglement-form one rather than merely equivalent to it. It is also why no entanglement
lower bound has to be tracked anywhere in this chapter.
\end{remark}

\begin{remark}[Relation to the abstract compression lemma]
\label{rem:compression-abstract}
\uses{thm:compression,lem:compressible-criterion,def:succinct,lem:value-lower-approx}
Theorem~\ref{thm:compression} instantiates the hypotheses of
Lemma~\ref{lem:compressible-criterion}: take $A$ to be the set of (descriptions of) normal
form verifier games with a value-$1$ PCC strategy and $B$ the set of those with $\valstar
\leq \frac12$, with the trivially accepting and the trivially rejecting game as the
distinguished elements; the complement of $B$ is semidecided by enumerating tensor-product
strategies, the $\valstar$ half of Lemma~\ref{lem:value-lower-approx}; and a description
$(\sampler, \decider, n)$ with $n$ in binary is a succinct description of the game
$\verifier_{2^n}$ played by a $\lambda$-bounded verifier. Completeness of compression is
the preservation of $A$; preservation of $B$ is the soundness clause of
Theorem~\ref{thm:compression} verbatim, hypothesis and conclusion being in the same value.
Taking $B$ in $\valstar$ throughout is what keeps Theorem~\ref{thm:almost-sync} out of the
pipeline altogether --- at the outer end its hypothesis is unavailable
(Remark~\ref{rem:almost-sync-hypothesis}) and at the inner end it is simply not needed
(Remark~\ref{rem:bipartite-route}); it is also the form in which the semidecision
procedure is available, the rounding argument for the $\synval$ half not being the same
one. Making this dictionary precise is a
formalization task of this section, and four obligations in it are open: the universe of
the criterion's strings --- if they are games, $\compr$ cannot be polynomial-time, a game
of $\verifier_{2^n}$ having $2^{\poly(n,\lambda)}$ questions, so they must be descriptions
and the interpretation of a description as a game must sit outside the time bound;
compression from a \emph{succinct} input rather than from a verifier read verbatim, as
Definition~\ref{def:succinct} and the polynomial-time convention of
Section~\ref{sec:conventions} require; preservation of $A$ and $B$ on inputs that are not
$\lambda$-bounded descriptions, a property the criterion demands of every string while
$\lambda$-boundedness is undecidable; and the choice of $\lambda$ along the recursion,
whose descriptions grow with the level. The entanglement form of~\cite{JNVWY20}
(Remark~\ref{rem:entanglement-form}) instantiates Lemma~\ref{lem:recursive-compression}
instead, with $f = \Ent(\cdot, \frac12)$; it is not used by the main theorem.
\end{remark}

\begin{remark}[The entanglement form of~\cite{JNVWY20}]
\label{rem:entanglement-form}
\uses{def:ent,thm:bvy,lem:recursive-compression}
The paper states the soundness of each transformation with an additional entanglement
clause, obtained from the same proofs by tracking the dimension of the extracted
strategies: $\Ent(\vint_n, 1 - \eps) \geq \max\{\Ent(\verifier_{2^n}, 1 -
\delta(\eps, n)),\ (1 - \delta(\eps, n))\, 2^{2^{\lambda n}}\}$ for introspection;
$\Ent(\vora_n, 1 - \eps) \geq \Ent(\verifier_n, 1 - \delta(\eps))$ for
oracularization; $\Ent(\vans_n, 1 - \eps) \geq \frac12 \Ent(\verifier_n, 1 - \delta(\eps, n))$
for answer reduction, the factor $\frac12$ being the Schmidt-rank doubling of the
symmetrization step, which the synchronous framework of this blueprint does not perform
(symmetry being definitional there) and therefore does not pay; Theorem~\ref{thm:bvy} for anchored repetition; and, composed,
$\Ent(\vcompr_n, \frac12) \geq \frac12 \max\{\Ent(\verifier_N, \frac12),\
2^{N^\lambda - 1}\}$ for compression. Fed into the variant of
Lemma~\ref{lem:recursive-compression} tolerating the factor $\frac12$ (harmless, since
the second term grows much faster than the levels of the recursion), with $f =
\Ent(\cdot, \frac12)$, this yields Theorem~\ref{thm:halting} with the stronger
conclusion $\Ent(\game_\machine, \frac12) = \infty$ for non-halting $\machine$. The
value form adopted here needs none of it; the entanglement clauses remain of interest for
the corollaries of Remark~\ref{rem:further}.

One proviso on the measure: the clauses above are transcribed from~\cite{JNVWY20}, where
the quantity is the least local dimension of a \emph{bipartite} tensor-product strategy,
whereas $\Ent$ of Definition~\ref{def:ent} is the least dimension of a \emph{synchronous}
one. Relating the two is the same transport as
Theorem~\ref{thm:almost-sync} and carries its hypothesis
(Remark~\ref{rem:almost-sync-hypothesis}), so the clauses should be read in the paper's
measure until that is done.
\end{remark}

\subsection{The halting reduction}
\label{sec:ps-halting}

The ingredients, with provenance to the companion repository's ledger nodes.

\begin{lemma}[The self-referential decider is well defined]
\label{lem:halt-construction}
\ledgernode{1.6.1}
\uses{def:normal-verifier,lem:compress-sampler-indep,lem:kleene}
There is a Turing machine $F$ and a decider $\dhalt = F$ carrying its own description,
with a timeout bound $B_{\dhalt}(n)$, such that $\dhalt$ is total within its timeout and
$\vhalt = (\mathtt{ComputeSampler}(\lambda), \dhalt)$ is a normal form verifier. The
self-reference is by direct construction rather than by the recursion theorem, the
input-length check precedes the $n$-step simulation of $\machine$, and the
halting-freeze convention makes the decider total.
\end{lemma}

\effortMedium

\begin{lemma}[A uniform bound on the verifier's parameter]
\label{lem:lambda-bound}
\ledgernode{1.6.4}
For all integers $C, C' \geq 1$, all $\lambda \geq 4 \max\{ (4C)^{8C},\ C \log C' \}$
and all $n \geq 2$,
\begin{equation*}
  C\, (C' \lambda n)^{C} \;\leq\; n^{\lambda} .
\end{equation*}
\end{lemma}

\effortEasy

\begin{lemma}[The parameter of the halting verifier]
\label{lem:lambda}
\ledgernode{1.6.5}
\uses{lem:halt-construction,lem:lambda-bound,def:lambda-bounded}
There is a $\lambda$, computable from the description of $\machine$ in polynomial time and
of size $\poly(\abs{\machine})$, such that $\vhalt$ is $\lambda$-bounded and
$\TIME_{\dhalt}(n), \TIME_{\shalt}(n) = \poly(n, \abs{\machine})$.
\end{lemma}

\effortMedium

\begin{lemma}[Values of the halting verifier]
\label{lem:dhalt-values}
\ledgernode{1.6.2, 1.6.3}
\uses{lem:halt-construction,lem:lambda,lem:compress-sampler-indep,def:pcc,def:q-value}
For every $n$:
\begin{enumerate}
\item if $\machine$ halts within $n$ steps then $\valstar(\vhalt_n) = 1$, witnessed by a
  value-$1$ PCC strategy;
\item if $\machine$ does not halt within $n$ steps then $\vhalt_n$ and $\vcompr_n$ are the
  same game up to an identification of the answer alphabets; in particular
  $\valstar(\vhalt_n) = \valstar(\vcompr_n)$, and $\vhalt_n$ has a value-$1$ PCC strategy
  if and only if $\vcompr_n$ does.
\end{enumerate}
\end{lemma}
\begin{proof}
Item 1 is the fixed-answer strategy: when $\machine$ halts within $n$ steps, $\dhalt$
accepts every answer pair on every question pair in the support, so any strategy has value
$1$ and the constant-$0$ strategy is PCC. For item 2, $\dhalt$ at index $n$ performs the
same question-length check as $\dcompr$ and then runs $\dcompr(n, x, y, a, b)$ verbatim,
with $\mathtt{TIMEOUT}$ rejecting in both, so the two deciders accept on exactly the same
inputs; and the two verifiers share the sampler $\shalt = \mathtt{ComputeSampler}(\lambda)$
(Lemma~\ref{lem:compress-sampler-indep}), so a strategy for either is a strategy for the
other with the same winning probability. The timeout counter of $\dhalt$ does not
interfere: its bound is chosen in Lemma~\ref{lem:lambda} so that the computations above
finish first.
\end{proof}

\begin{remark}[Item 2 is a transfer, not a bound]
\label{rem:dhalt-transfer}
\uses{lem:dhalt-values,thm:halting,thm:compression}
Item 2 is deliberately not ``$\valstar(\vhalt_n) \leq \frac12$''. That per-level bound is
false: if $\machine$ halts for the first time at step $T$, then for every $n < T$ it does
not halt within $n$ steps, yet $\valstar(\vhalt_n) = 1$ --- item 1 gives value $1$ at every
level $\geq T$, and item 2 together with the completeness clause of
Theorem~\ref{thm:compression} carries it \emph{downward} from level $2^n$ to level $n$,
which is exactly how the completeness half of Theorem~\ref{thm:halting} is proved. The
absolute bound needs $\machine$ never to halt, and is then the conclusion of the recursion
in Theorem~\ref{thm:halting}, not an ingredient of it. Recorded here because an earlier
revision of this chapter did state item 2 as the bound, and because the ledger node
`1.6.3' states it that way too; the formalization should follow the paper.
\end{remark}

\effortHard

\begin{theorem}[Halting reduction]
\label{thm:halting}
\ledgernode{1.6, 1.6.6}
\uses{thm:compression,rem:compression-abstract,lem:compressible-criterion,lem:value-lower-approx,lem:kleene,def:mipstar,def:game-description,lem:sync-le-valstar,lem:dhalt-values,lem:lambda}
There is a computable (indeed polynomial-time) map $\machine \mapsto
\game_{\machine}$ from Turing machines to synchronous game descriptions, with games
given by normal form verifiers satisfying the efficiency requirements of
Definition~\ref{def:mipstar}, such that:
\begin{enumerate}
\item if $\machine$ halts on the empty input then $\game_\machine$ has a value-$1$ PCC
  strategy, so $\synval(\game_\machine) = \valstar(\game_\machine) = 1$;
\item if $\machine$ does not halt on the empty input then
  $\valstar(\game_\machine) \leq \tfrac12$, and hence also
  $\synval(\game_\machine) \leq \tfrac12$ (Lemma~\ref{lem:sync-le-valstar}).
\end{enumerate}
\end{theorem}

The proof is Lemma~\ref{lem:compressible-criterion} applied to the instantiation of
Remark~\ref{rem:compression-abstract}; the semidecidability hypothesis of the lemma is the
$\valstar$ half of Lemma~\ref{lem:value-lower-approx}. The obligations of that
instantiation that are open are inherited here: the four listed in
Remark~\ref{rem:compression-abstract}. The synchronization invariant that an earlier
revision also required of the games produced inside $\Compress$ is no longer among them
(Remark~\ref{rem:bipartite-route}).

\begin{remark}[Identical measurement operators]
\label{rem:identical-operators}
\ledgernode{1.6.7}
\uses{thm:halting,def:sync-strategy,def:pcc,thm:separation}
One property the synchronous framework gets for nothing is worth recording. The value-$1$
witness of item 1 is a PCC, hence synchronous, strategy: a \emph{single} family of
measurements used by both players. In the bipartite framework of~\cite{JNVWY20} that the
two players measure identically is an extra property, established by tracking it through
oracularization, answer reduction, detyping and repetition, and the verification
campaign~\cite{Audit26} reports that the chain is delicate, since the completeness clause
of compression carries only PCC and the property has to be read off the last application
rather than carried by the induction. Here it holds by definition
(Definition~\ref{def:sync-strategy}), so no such chain is needed --- which is also what
makes the commuting completeness of Theorem~\ref{thm:separation} a question about one
transformation at a time rather than about a provenance argument. The entanglement form of the statement, with the
stronger conclusion $\Ent(\game_\machine, \frac12) = \infty$ for non-halting
$\machine$, is discussed in Remark~\ref{rem:entanglement-form}.
\end{remark}
